(1) 以热端放出的热率 \(J_q\)、冷端得到的热率 \(J_q-P\) 计算热库熵变化，
\[
\dot S_{\rm i}=\frac{J_q-P}{T_c}-\frac{J_q}{T_h}
=\frac{J_q\Delta T}{T^2}-\frac{P}{T}+O((\Delta T/T)^3)
=I X_1+J_qX_2
\]
（这里电压、电流、热流均为小温差的一阶量；余项只表示超出二阶熵产生精度的修正）。Onsager 对称性使
\[
\dot S_{\rm i}=L_{11}X_1^2+2L_{12}X_1X_2+L_{22}X_2^2
=L_{11}\left(X_1+\frac{L_{12}}{L_{11}}X_2\right)^2+
\frac{\det L}{L_{11}}X_2^2 .
\]
对任意力非负的充要条件是 \(L_{11}\ge0,L_{22}\ge0,\det L\ge0\)；本题严格正定的假设使两项系数均正。由于 \(L_{12}>0\)，
\[
\boxed{0<q=\frac{L_{12}}{\sqrt{L_{11}L_{22}}}<1},\qquad
\frac{\det L}{L_{11}}=L_{22}(1-q^2)>0.
\]
第一平方项是电力偏离开路补偿值所引起的耗散，第二项是无法由电负载消除的漏热耗散。它们分别在 \(X_1=-L_{12}X_2/L_{11}\) 和 \(\det L\to0^+\) 时趋于零；不能只检验对角元正就忽略交叉耦合的正定约束。

(2) 代入负载参数 \(x\)，直接得到
\[
\boxed{I=L_{12}X_2(1-x)},\qquad
\boxed{J_q=L_{22}X_2(1-q^2x)},\qquad
\boxed{P=T\frac{L_{12}^2}{L_{11}}X_2^2x(1-x)} .
\]
由(1)的完成平方式，
\[
\boxed{\dot S_{\rm i}=L_{22}X_2^2\left[(1-q^2)+q^2(1-x)^2\right]
=L_{22}X_2^2[1-q^2(2x-x^2)]>0}.
\]
在 \(x=0\) 时 \(V=0\) 为短路，\(I=L_{12}X_2\) 而 \(P=0\)；在 \(x=1\) 时 \(I=0\) 为开路，\(J_q=L_{22}X_2(1-q^2)>0\) 而 \(P=0\)。对 \(0<x<1\)，\(P>0\) 且 \(J_q\ge L_{22}X_2(1-q^2)>0\)。所以
\[
\boxed{\frac{\eta}{\eta_{\mathrm C}^{(1)}}
=\frac{P/J_q}{TX_2}=\frac{q^2x(1-x)}{1-q^2x}}.
\]
由能量与熵账本恒等式
\[
\eta_{\mathrm C}^{(1)}J_q-P=T\dot S_{\rm i}\ge0
\]
可知输出功率不超过该精度下的卡诺可用热功。注意 \(J_q\) 与 \(P\) 的阶数分别为 \(\Delta T\) 与 \((\Delta T)^2\)，故上述效率只保留首个非零阶。

(3) \(P\) 的负载因子 \(x(1-x)=\tfrac14-(x-\tfrac12)^2\)，故
\[
\boxed{x_P=\frac12},\qquad
\boxed{P_{\max}=\frac{TL_{12}^2X_2^2}{4L_{11}}},\qquad
\boxed{\frac{\eta(P_{\max})}{\eta_{\mathrm C}^{(1)}}=
\frac{q^2}{2(2-q^2)}\le\frac12}.
\]
设 \(z=q^2\in(0,1)\)。对 \(f(x)=zx(1-x)/(1-zx)\) 求导，
\[
f'(x)=\frac{z(1-2x+zx^2)}{(1-zx)^2}.
\]
分子在区间 \([0,1]\) 内只有一个零点
\[
\boxed{x_\eta=\frac{1-\sqrt{1-z}}{z}
=\frac1{1+\sqrt{1-q^2}}},
\]
且导数从正变负，故它是唯一全局效率极大点。令 \(s=\sqrt{1-q^2}\)，代入得
\[
\boxed{\frac{\eta_{\max}}{\eta_{\mathrm C}^{(1)}}=
\frac{1-s}{1+s}
=\frac{\sqrt{1+ZT}-1}{\sqrt{1+ZT}+1}},\qquad
ZT=\frac{q^2}{1-q^2}.
\]
弱耦合 \(q\to0\) 时最大效率趋零，且 \(P_{\max}\propto L_{12}^2\)；严格说若同时变化其他 \(L_{ij}\)，绝对功率极限还取决于其尺度。紧耦合 \(q\to1^-\) 时效率上界趋向 \(\eta_{\mathrm C}^{(1)}\)，但 \(x_\eta\to1\)、\(I\to0\)，在固定 \(L_{11},L_{22},X_2\) 下最优效率点的 \(P\to0\)。相反最大功率点仍有不可逆熵产生，其效率最多为线性响应卡诺效率的一半。

(4) 用 \(V=-TX_1,\ \Delta T=T^2X_2\) 逐项代回给定本构关系，
\[
I=GTX_1+GST^2X_2,\qquad
J_q=GST^2X_1+(GS^2T^3+KT^2)X_2.
\]
因此
\[
\boxed{L_{11}=GT,\quad L_{12}=GST^2,\quad
L_{22}=GS^2T^3+KT^2},\qquad
\boxed{\det L=GKT^3>0}.
\]
两交叉系数相等正是无磁场的 Onsager 互易性；Peltier 热系数 \(\Pi=ST\) 保证这一相等性。于是
\[
q^2=\frac{GS^2T}{GS^2T+K},\qquad
\boxed{ZT=\frac{q^2}{1-q^2}=\frac{GS^2T}{K}}.
\]
负载电压 \(V=xS\Delta T\)、电流 \(I=GS\Delta T(1-x)\)，故
\[
\frac{R_L}{1/G}=\frac{x}{1-x},\qquad
\boxed{R_{L,P}=\frac1G},\qquad
\boxed{R_{L,\eta}=\frac{\sqrt{1+ZT}}G}.
\]
直接代入物理流可得
\[
\boxed{\dot S_{\rm i}=-\frac{IV}{T}+\frac{J_q\Delta T}{T^2}
=\frac{G(S\Delta T-V)^2}{T}+\frac{K(\Delta T)^2}{T^2}\ge0}.
\]
开路 \(V=S\Delta T\) 时 \(I=0\)、\(J_q=K\Delta T\)，证明 \(K\) 是无法以电负载消除的漏热系数。\(K\to0^+\) 时 \(ZT\to\infty\) 且 \(R_{L,\eta}\to\infty\)；效率逼近线性响应卡诺值的同时最优效率负载电流与输出功率趋零。量纲检查：\(G\) 为 \(\mathrm{A/V}\)、\(S\) 为 \(\mathrm{V/K}\)、\(K\) 为 \(\mathrm{W/K}\)，因此 \(GS^2T/K\) 无量纲，\(\dot S_{\rm i}\) 两项均为 \(\mathrm{W/K}\)。
